\section{Perangkat dan Lingkungan Eksperimen}

Seluruh algoritma dijalankan pada perangkat dan lingkungan eksperimen yang identik sesuai variabel kontrol. 
Subbab ini menjelaskan perangkat keras dan perangkat lunak yang digunakan, serta prosedur pengukuran memori dan waktu komputasi pada aktivitas \textit{Demonstration} dan \textit{Evaluation}.

Perangkat target adalah Raspberry Pi 5 varian RAM\label{first:ram} 8 GB LPDDR4X-4267 dengan \textit{System-on-Chip Broadcom BCM2717}, yang memiliki \textit{CPU\label{first:cpu} quad-core ARM Cortex-A76 2,4 GHz}, \textit{cache} L2 512 KB per inti, dan \textit{cache} L3 bersama 2 MB \parencite{raspberry_pi_ltd_raspberry_2026}. 
Perangkat menggunakan penyimpanan microSD dan Raspberry Pi OS 64-bit berbasis Debian. 
Konfigurasi ini digunakan untuk merepresentasikan perangkat \textit{edge-IoT} dengan sumber daya komputasi terbatas dibandingkan \textit{server} konvensional.

CoDA-EDA dan seluruh \textit{baseline} diimplementasikan menggunakan Python. 
Pemilihan ini konsisten dengan ekosistem implementasi EDA yang tersedia; BOA dan berbagai EDA multivariat, misalnya, telah tersedia dalam paket Python seperti EDAspy \parencite{soloviev_edaspy_2024}, sehingga implementasi \textit{baseline} dapat mengacu pada implementasi referensi yang tersedia.

Jejak memori diukur menggunakan dua instrumen yang saling melengkapi. 
Metrik utama, $MemPeak_{alloc}$, merupakan puncak alokasi memori struktur data algoritma yang diukur menggunakan $tracemalloc$. 
Metrik kedua, $MemPeak_{sys}$, diperoleh dari VmHWM\label{first:vmhwm} pada $/proc/[pid]/status$ untuk memberikan konteks penggunaan memori pada tingkat sistem operasi. 
Pemisahan ini diperlukan karena VmHWM mencakup \textit{overhead} interpreter dan pustaka; pada perangkat target, nilai awalnya sekitar 31 MB, sedangkan perbedaan struktur data antaralgoritma dapat berada pada skala \textit{kilobyte} hingga beberapa \textit{megabyte}. 
Dengan demikian, penggunaan VmHWM saja dapat menyamarkan perbedaan memori algoritmik.

Kedua metrik dicatat sebagai nilai puncak karena kelayakan eksekusi pada perangkat \textit{edge} ditentukan oleh kebutuhan memori maksimum selama algoritma berjalan. 
Setiap kombinasi algoritma dan dataset dieksekusi dalam proses sistem operasi yang terpisah karena VmHWM bersifat \textit{high-water mark} dan tidak menurun selama umur proses. 
Pendekatan ini juga mencegah pengukuran sebelumnya memengaruhi konfigurasi berikutnya. 
Eksekusi dilakukan satu per satu tanpa algoritma lain yang berjalan secara bersamaan untuk mengurangi interferensi penggunaan sumber daya.

Waktu komputasi diukur menggunakan \textit{wall-clock time} pada perangkat target dan dibedakan menjadi waktu pembelajaran awal (\textit{initial training}) serta waktu satu siklus pembaruan inkremental. 
Untuk mengurangi variasi akibat \textit{dynamic frequency scaling} dan \textit{thermal throttling}, \textit{CPU governor} ditetapkan pada mode \textit{performance} selama pengujian dan suhu perangkat dipantau agar tetap berada dalam rentang operasi yang ditetapkan produsen.

\begin{longtable}{|>{\raggedright\arraybackslash}p{3cm}|>{\raggedright\arraybackslash}p{9cm}|}
\caption{Spesifikasi Perangkat dan Lingkungan Eksperimen}
\label{tab:spesifikasi-perangkat} \\
\hline
\textbf{Komponen} & \textbf{Spesifikasi/Prosedur} \\
\hline
\endfirsthead

\multicolumn{2}{c}{{\bfseries \tablename\ \thetable{} -- lanjutan dari halaman sebelumnya}} \\
\hline
\textbf{Komponen} & \textbf{Spesifikasi/Prosedur} \\
\hline
\endhead

\hline
\multicolumn{2}{r}{Lanjut ke halaman berikutnya} \\
\endfoot

\hline
\endlastfoot

% ==================== DATA BARIS ====================
Perangkat & Raspberry Pi 5 (Raspberry Pi Ltd., 2026) \\
\hline
SoC\label{first:soc} & Broadcom BCM2712, quad-core ARM\label{first:arm} Cortex-A76 @ 2{,}4GHz, cache L2 512KB/core, cache L3 bersama 2MB \\
\hline
RAM & 8GB LPDDR4X-4267 SDRAM\label{first:sdram} \\
\hline
Penyimpanan & microSD (mode SDR104); PCIe\label{first:pcie} 2.0 x1 tersedia untuk SSD\label{first:ssd} via M.2 HAT\label{first:hat} terpisah (tidak dipakai sebagai konfigurasi baku) \\
\hline
Suhu Operasi & $0^\circ$C--$70^\circ$C (spesifikasi produsen) \\
\hline
Sistem Operasi & Raspberry Pi OS\label{first:os} 64-bit (berbasis Debian) \\
\hline
Bahasa Pemrograman & Python \\
\hline
Pengukuran Memori & Dual: puncak alokasi \texttt{tracemalloc} (metrik utama) + VmHWM pada /proc/[pid]/status (konteks sistem); satu proses terpisah per konfigurasi \\
\hline
Pengukuran Waktu & Wall-clock time langsung pada perangkat target; CPU governor mode \texttt{performance}; suhu dipantau terhadap rentang operasi resmi \\
\hline

\end{longtable}